\documentclass[11pt]{article}

\usepackage[utf8]{inputenc}
\usepackage[margin=1in]{geometry}
\usepackage{times}
\usepackage{amsmath}
\usepackage{amssymb}
\usepackage{booktabs}
\usepackage{listings}
\usepackage{xcolor}
\usepackage{hyperref}
\usepackage{enumitem}
\usepackage{caption}

\hypersetup{colorlinks=true, linkcolor=blue, citecolor=blue, urlcolor=blue}

\lstset{
  basicstyle=\ttfamily\footnotesize,
  breaklines=true,
  frame=single,
  columns=fullflexible,
  keepspaces=true,
  showstringspaces=false,
  backgroundcolor=\color{gray!8},
}

\title{\textbf{Standardizing Data Staleness Telemetry:\\
A Vendor-Neutral OpenTelemetry Semantic Convention\\
for Data Freshness}}

\author{Anirudh Rajmohan}
\date{\today}

\begin{document}
\maketitle

\begin{abstract}
Data freshness---how recently a dataset reflects the real world---is a
first-class reliability concern for analytics, machine learning, and
operational decision systems. Yet despite decades of theory on information
freshness (the \emph{Age of Information}, AoI) and a thriving market of
data-observability tools, there is no portable, vendor-neutral way to
\emph{emit} freshness as telemetry. OpenTelemetry (OTEL), now the de facto
standard for observability signals, defines semantic conventions for database
and messaging operations but none for data staleness; freshness today is
re-implemented, incompatibly, inside each observability vendor and each data
platform. We address this gap with three contributions. First, we propose a
compact OpenTelemetry semantic convention for data staleness---a small set of
metrics (\texttt{data.staleness.age}, \texttt{lag}, \texttt{last\_update},
\texttt{records.behind}, and SLA threshold/breach signals) and attributes that
unify freshness across SQL warehouses, streaming systems, batch pipelines,
caches, NoSQL and lakehouse stores, and search/vector indexes under one
vocabulary, grounded in the AoI model, including a \emph{differential}
(source-relative) form for indexes and replicas that trail their upstream.
Second, we provide a
reference implementation: a Python SDK that turns per-source \emph{probes} into
standardized metrics, a zero-configuration OpenTelemetry Collector receiver that
scrapes SQL, streaming, object, and version-drift sources with no application
code, and a Collector processor that centrally derives age and evaluates
freshness SLAs. Third, we evaluate the
implementation's correctness and overhead, measuring per-source emission cost
of \(22\)--\(39\,\mu s\) in the SDK and sub-microsecond per-point processing in
the Collector. We further validate the implementation end-to-end against real
backends---PostgreSQL, Kafka, Redis, Kinesis, and a schema registry---on a
cloud instance, asserting that emitted metrics are numerically correct (age
equal to a known injected offset, exact consumer lag) and that failures are
made visible rather than fabricated, not merely that metrics appear. Our aim is not a new freshness metric but the
\emph{consolidation} of a well-studied one into the dominant observability
framework, so that freshness becomes as portable and queryable as latency or
error rate. The artifacts are released as open source under Apache-2.0.
\end{abstract}

\section{Introduction}
Modern organizations make decisions on derived data: dashboards, feature
stores, recommendation indexes, fraud models, and financial reports all consume
data that has traveled through ingestion, transformation, and serving layers.
When that pipeline stalls, the data does not become \emph{wrong} so much as
\emph{old}: a dashboard silently shows yesterday's numbers, a model scores on
features that no longer reflect reality. This failure mode---\emph{staleness}---
is distinct from the failures that traditional observability captures well.
A pipeline can have zero errors and low latency on every individual operation
and still serve hours-old data because an upstream producer stopped.

Two largely separate communities have studied this. In networked and
cyber-physical systems, the \emph{Age of Information} (AoI) line of work
formalizes freshness as the time since the most recently received update was
generated at its source~\cite{kaul2012,yates2021survey}, with a rich body of
analysis and variants such as Peak AoI and Age of
Synchronization~\cite{yates2021survey}. In data engineering, a generation of
\emph{data-observability} platforms~\cite{montecarlo,metaplane,sifflet} and
tools such as dbt \texttt{source freshness}~\cite{dbtfreshness} and Kafka lag
monitors like Burrow~\cite{burrow} measure freshness operationally. Both agree
on the core quantity. Neither produces it as portable telemetry.

The consequence is fragmentation. Each data-observability vendor models and
stores freshness differently; freshness computed by dbt cannot be compared,
in one query, with freshness of a Kafka topic or a Redis cache; and alerting on
freshness is locked to whichever proprietary system computed it. Meanwhile
OpenTelemetry (OTEL) has become the cross-vendor standard for emitting metrics,
traces, and logs, with \emph{semantic conventions} that make signals
comparable across systems~\cite{otelsemconv}. OTEL defines conventions for
database client metrics and messaging/Kafka metrics, but these capture
operation \emph{duration} and message/row \emph{counts}---not how old the data
is. As of the 1.42 conventions and the published 2026 semantic-conventions
roadmap, there is no data-freshness metric, and not even a standardized
consumer-lag metric~\cite{otelmessaging,otelroadmap}.

This paper closes that gap. We argue that freshness deserves the same
treatment latency received: a small, agreed-upon convention plus reference
instrumentation, so any producer can emit it and any backend can consume it.
Concretely, we contribute:
\begin{enumerate}[nosep]
  \item \textbf{A semantic convention for data staleness} (Section~\ref{sec:conv}):
  metric names, units, instrument types, and attributes that express AoI-style
  freshness uniformly across SQL/warehouse, streaming, batch-pipeline, and
  cache/object-store sources.
  \item \textbf{A reference implementation} (Sections~\ref{sec:design}--\ref{sec:impl}):
  an Apache-2.0 Python SDK (\texttt{otel-staleness}) built on
  \texttt{opentelemetry-python}, and a Go OpenTelemetry Collector processor
  (\texttt{datastaleness}) that derives age and evaluates SLAs centrally.
  \item \textbf{An evaluation} (Section~\ref{sec:eval}) of correctness and
  runtime overhead, with reproducible microbenchmarks.
\end{enumerate}
We are explicit about novelty: the \emph{metric} of staleness is not new---it
is AoI, and it is what data-observability vendors already compute. Our
contribution is \emph{standardization and consolidation}: expressing that known
quantity once, in OpenTelemetry, with a portable implementation.

\section{Background and Related Work}
\label{sec:related}

\paragraph{Age of Information.}
AoI was introduced to capture freshness in status-update
systems~\cite{kaul2012} and has since become a canonical metric for
information freshness in communication, IoT, and control
systems~\cite{yates2021survey}. For a source whose freshest delivered update
was generated at time \(t_e\), the age at wall-clock time \(t\) is
\(\Delta(t) = t - t_e\): a sawtooth that rises linearly until a newer update
arrives and resets it. Our \texttt{data.staleness.age} is exactly this
quantity, applied to data systems rather than wireless links. AoI theory also
motivates our distinction between \emph{age} (which grows during a stall) and
one-shot \emph{lag} (the delay a record incurs traversing the pipeline).

\paragraph{Data freshness in data management.}
The data-management community studied freshness and timeliness well before
``data observability'' as a term, e.g.\ Peralta's survey of data freshness and
accuracy in integration systems~\cite{peralta2006}. Freshness is now treated as
a pillar of data observability alongside volume, schema, distribution, and
lineage~\cite{montecarlo}. Practitioner definitions converge on the lag between
when an event occurred and when it is queryable~\cite{metaplane,sifflet}.

\paragraph{Operational tooling.}
dbt computes source freshness from a \texttt{loaded\_at} column and warns/errors
against configured thresholds~\cite{dbtfreshness}; Kafka ecosystems track
consumer lag with tools such as Burrow~\cite{burrow}; commercial platforms
(Monte Carlo, Metaplane, Sifflet) provide freshness monitoring as a managed
service~\cite{montecarlo,metaplane,sifflet}. Each is valuable and each is a
silo: the freshness signal lives inside the tool that produced it.

\paragraph{OpenTelemetry.}
OTEL standardizes telemetry emission and, via semantic conventions, the
\emph{meaning} of that telemetry~\cite{otelsemconv}. Conventions exist for
database and messaging metrics~\cite{otelmessaging}, and the data model includes
a low-level Prometheus-style ``staleness marker'' for missing
timeseries---an unrelated concept to data freshness. No convention emits data
age. Our work proposes one and is structured to serve as the basis for an
OpenTelemetry Enhancement Proposal (OTEP).

\paragraph{Positioning.}
Table~\ref{tab:related} summarizes how existing approaches relate to the three
properties we care about: a formal freshness model, operational measurement
across heterogeneous sources, and portable/vendor-neutral emission.

\begin{table}[t]
\centering
\caption{Where existing approaches stand. ``Portable emission'' means the
freshness signal is produced as open, cross-vendor telemetry.}
\label{tab:related}
\begin{tabular}{lccc}
\toprule
Approach & Formal model & Multi-source & Portable emission \\
\midrule
AoI theory~\cite{kaul2012,yates2021survey} & \checkmark & partial & --- \\
dbt source freshness~\cite{dbtfreshness} & partial & --- (batch only) & --- \\
Kafka lag tools~\cite{burrow} & partial & --- (streaming only) & --- \\
Data-observability SaaS~\cite{montecarlo,metaplane} & partial & \checkmark & --- \\
OTEL DB/messaging semconv~\cite{otelmessaging} & --- & partial & \checkmark \\
\textbf{This work} & \checkmark & \checkmark & \checkmark \\
\bottomrule
\end{tabular}
\end{table}

\section{The Data Staleness Convention}
\label{sec:conv}

\subsection{Model}
Let a \emph{logical source} be any dataset whose freshness we care about: a
table, a topic, a dbt model, a cache namespace, or an object-store prefix. For
the freshest record of a source with event time \(t_e\), observed/processed at
\(t_p\), with current wall-clock time \(t_{now}\):
\begin{align}
\text{lag} &= t_p - t_e, \\
\text{age} &= t_{now} - t_e \;\;(\geq \text{lag}).
\end{align}
\texttt{lag} is the delay a record incurred moving through the pipeline; it is
fixed per record. \texttt{age} keeps increasing until a newer record arrives.
If a source stops producing, lag of the last record is constant while age rises
without bound---precisely the condition freshness monitoring exists to detect.

\subsection{Metrics}
All durations use the UCUM unit \texttt{s}; timestamps are Unix epoch seconds.
Point-in-time values are reported as asynchronous (observable) gauges; counts of
events use monotonic counters. Table~\ref{tab:metrics} lists the convention.

\begin{table}[t]
\centering
\caption{Proposed \texttt{data.staleness.*} metrics.}
\label{tab:metrics}
\footnotesize
\begin{tabular}{lllp{5.4cm}}
\toprule
Metric & Instrument & Unit & Meaning \\
\midrule
\texttt{data.staleness.age} & Gauge & \texttt{s} & \(t_{now}-t_e\); primary AoI metric. \\
\texttt{data.staleness.lag} & Gauge & \texttt{s} & \(t_p-t_e\) of most recent record. \\
\texttt{data.staleness.last\_update.timestamp} & Gauge & \texttt{s} & Unix time of last update. \\
\texttt{data.staleness.records.behind} & Gauge & \texttt{\{record\}} & Positional backlog (e.g.\ Kafka lag). \\
\texttt{data.staleness.sla.threshold} & Gauge & \texttt{s} & Max acceptable age (SLO target). \\
\texttt{data.staleness.sla.breached} & Gauge & \texttt{1} & \(1\) if age \(>\) threshold else \(0\). \\
\texttt{data.staleness.sla.breaches} & Counter & \texttt{\{breach\}} & Transitions into breach. \\
\bottomrule
\end{tabular}
\end{table}

\subsection{Attributes}
A small attribute set makes points comparable and sliceable:
\texttt{data.source.system} (e.g.\ \texttt{postgresql}, \texttt{kafka},
\texttt{dbt}, \texttt{redis}, \texttt{s3}; reusing \texttt{db.system}/
\texttt{messaging.system} values where they exist),
\texttt{data.source.name} (table/topic/model/prefix),
\texttt{data.source.namespace} (schema/consumer group/bucket),
\texttt{data.staleness.method} (how the value was derived; see below),
\texttt{data.staleness.partition}, and \texttt{data.pipeline.stage}
(\texttt{ingest}/\texttt{transform}/\texttt{serve}).

The \texttt{method} attribute records the measurement recipe and makes
heterogeneous sources self-describing: \texttt{max\_timestamp} (SQL
\texttt{MAX(updated\_at)}), \texttt{watermark}, \texttt{consumer\_lag}
(streaming), \texttt{run\_completion} (batch jobs), \texttt{object\_mtime}
(object storage), \texttt{ttl\_age} (caches), and \texttt{heartbeat}.

\subsection{Per-source mapping}
The convention is useful only if each source type has an unambiguous recipe.
SQL/warehouse sources set \(age = t_{now} - \texttt{MAX}(\textit{ts column})\).
Streaming sources use the event timestamp of the last consumed record for age
and lag, and offset lag for \texttt{records.behind}, reported per partition.
Batch/dbt/Airflow sources set \(age = t_{now} - \textit{last successful run
end}\). Caches and object stores use last write / last-modified time. These
mirror what existing tools already compute internally; the convention's value is
that they now emit \emph{the same metric}. Because the model rests only on a
source event time, it extends to further families with no new metrics---NoSQL
and graph stores (e.g.\ Cassandra \texttt{WRITETIME()}), time-series databases,
lakehouse table formats (Iceberg/Delta/Hudi snapshot commit time), and external
feeds---which become new \texttt{data.source.system} / \texttt{method} values
rather than new instruments.

\paragraph{Differential (source-relative) freshness.}
Some systems' freshness is best expressed against an \emph{upstream} rather than
against \(t_{now}\): a search or vector index that trails its source corpus, or
a replica/CDC target that trails its primary. We capture this with the existing
\texttt{lag} metric, setting the record's event time to the upstream
write/commit time and naming the upstream in a \texttt{data.staleness.relative\_to}
attribute (methods \texttt{index\_lag}, \texttt{replication\_lag}). Using source
time on both sides makes the measurement robust to clock skew between the two
systems. Absolute (\texttt{age} vs \(t_{now}\)) and differential
(\texttt{lag} vs upstream) freshness are complementary and may be emitted
together; index staleness in retrieval-augmented generation pipelines is a
prominent case where the index can be behind its corpus while also being old in
absolute terms.

\section{System Design}
\label{sec:design}
The reference implementation has two cooperating halves
(described below): a client-side SDK that knows how to
\emph{measure} each source, and a Collector-side processor that can
\emph{derive and evaluate} freshness centrally.

\paragraph{Architecture (textual).}
\textit{Producers} expose freshness via \emph{probes}. A probe reports the event
time of the freshest record for one logical source. The \texttt{StalenessMonitor}
registers the standardized observable gauges once and, on each collection,
pulls readings from all probes, computes age, and emits observations. Telemetry
flows over OTLP to an OpenTelemetry Collector, where the \texttt{datastaleness}
processor can (a) derive \texttt{data.staleness.age} for producers that emit
only \texttt{last\_update.timestamp}, and (b) attach
\texttt{data.staleness.sla.threshold}/\texttt{breached} according to operator
policy. Any OTEL-compatible backend (Prometheus, OTLP store, vendor) then stores
and queries the result.

\paragraph{Why two halves.}
Some producers can compute age themselves (a service that knows its event
times); others are best instrumented with the bare minimum---``here is the
timestamp of my newest row''---and should not embed clock logic or SLA policy.
Centralizing derivation and SLA evaluation in the Collector keeps thresholds in
one operational place and lets the same age computation apply uniformly, which
also avoids clock-skew logic sprawling across services.

\paragraph{Design principles.}
(1) \emph{Minimum producer burden}: a probe is a plain callable returning a
timestamp. (2) \emph{Fault isolation}: a failing probe must never break
collection of the others. (3) \emph{No new transport}: everything rides
standard OTEL metrics. (4) \emph{Policy where it belongs}: SLAs are operator
configuration in the Collector, not hard-coded in applications.

\section{Implementation}
\label{sec:impl}

\subsection{Python SDK}
\texttt{otel-staleness} is built on \texttt{opentelemetry-api/sdk}. The core is
a \texttt{FreshnessReading} dataclass and a \texttt{StalenessMonitor} that
registers the convention's instruments via observable-gauge callbacks. Probes
ship for the four core families (SQL/warehouse, streaming, batch, cache/object
store) plus two \emph{differential} probes---\texttt{IndexFreshnessProbe} and
\texttt{ReplicationFreshnessProbe}---that report how far a search/vector index or
a replica trails its upstream via \texttt{lag} and \texttt{relative\_to}:

\begin{lstlisting}[language=Python]
from otel_staleness import StalenessMonitor, conventions as sc
from otel_staleness.probes import SQLFreshnessProbe

monitor = StalenessMonitor()
monitor.add_probe(SQLFreshnessProbe(
    fetch_max_epoch=lambda: latest_updated_at(),  # Unix s or datetime
    source_name="orders", system=sc.System.POSTGRESQL,
    namespace="public", sla_threshold_seconds=300))
monitor.start()
\end{lstlisting}

Every probe accepts a plain callable, so probes are unit-testable without a live
backend; \texttt{SQLFreshnessProbe} additionally offers a \texttt{from\_sqlalchemy}
constructor. The monitor guards each probe call so a raised exception is
swallowed and the remaining probes still report---design principle (2).

\subsection{Collector processor}
\texttt{datastaleness} is a Go metrics processor (Collector API
\texttt{v0.110.0}, pdata \texttt{v1.16.0}). On each batch it: scans for existing
\texttt{age} points; optionally derives \texttt{age} from
\texttt{last\_update.timestamp} as \(\max(0, t_{now}-ts)\); and, for each age
point, looks up the first matching SLA rule (selectors on
\texttt{data.source.system}/\texttt{name}, with a default threshold) to emit
\texttt{threshold} and \texttt{breached} points. Configuration:

\begin{lstlisting}
processors:
  datastaleness:
    compute_age_from_last_update: true
    evaluate_sla: true
    default_threshold: 5m
    slas:
      - source_system: kafka
        threshold: 30s
      - source_name: orders
        threshold: 1m
\end{lstlisting}

\section{Evaluation}
\label{sec:eval}
We evaluate (i) functional correctness and (ii) runtime overhead. Experiments
ran in a containerized Linux environment (Python 3.10, Go 1.25, 2 vCPU).

\subsection{Correctness}
Both components ship with unit tests asserting the convention's semantics:
age derived from last-update time; explicit age preferred when given; age
clamped to be non-negative under clock skew; SLA breach edges; per-partition
Kafka lag; method/attribute tagging; and probe fault isolation. The Python
suite (17 tests) and Go suite (7 tests) pass. Representative invariants checked:
\(age = t_{now}-t_e\); \(age \geq 0\); and \(\texttt{breached}=\mathbb{1}[age >
\textit{threshold}]\).

\subsection{Overhead: SDK}
We measured the cost of one collection cycle (callback execution plus in-memory
export) as a function of the number of monitored sources, averaged over 20
iterations. Table~\ref{tab:sdkbench} shows the SDK adds tens of microseconds per
source. Even at \(10{,}000\) sources a full collection completes in
\(\sim\!385\,\text{ms}\); at a typical 15--60\,s metric export interval this is
negligible. Per-source cost is dominated by attribute-set construction and
observation allocation, not by the freshness arithmetic.

\begin{table}[t]
\centering
\caption{Python SDK: collection + in-memory export cost.}
\label{tab:sdkbench}
\begin{tabular}{rrr}
\toprule
Sources & Cycle time (ms) & Per source (\(\mu s\)) \\
\midrule
100 & 2.16 & 21.7 \\
1{,}000 & 25.7 & 25.7 \\
10{,}000 & 385.3 & 38.5 \\
\bottomrule
\end{tabular}
\end{table}

\subsection{Overhead: Collector processor}
We benchmarked \texttt{processMetrics} on a batch of \(1{,}000\)
\texttt{last\_update.timestamp} points with both age derivation and SLA
evaluation enabled (Go \texttt{testing} benchmark). It runs in
\(\sim\!0.93\,\text{ms}\) per batch (\(\approx 0.93\,\mu s\) per point) with
\(\sim\!1.5\,\text{MB}\) allocated per batch. Throughput therefore exceeds one
million points per second per core, far above the point rates produced by
freshness telemetry, which is emitted per source per export interval rather than
per request.

\subsection{Discussion of results}
The measured overheads confirm the design intent: freshness is a low-cardinality,
low-frequency signal (one set of points per source per export), so even naive
implementations are comfortably cheap. The interesting cost in practice is
cardinality---number of distinct \texttt{(system, name, partition)} tuples---
which operators already manage for other metrics and which the attribute set is
deliberately kept small to bound.

\section{Limitations and Future Work}
\label{sec:limits}
Our evaluation establishes correctness and overhead, not end-to-end detection
quality on production workloads; a natural next step is a field study comparing
breach-detection latency against an existing data-observability tool on the same
pipelines. Clock skew between producer and Collector affects derived age; we
clamp to non-negative but do not yet correct skew (e.g.\ via NTP offset
attributes). The convention intentionally omits richer freshness variants from
AoI theory (Peak AoI, Age of Incorrect Information); these could be added as
optional metrics. Finally, the names here are a Development-stage proposal: the
intended path to adoption is an OTEP and community review, during which names may
change. Additional first-party probes (BigQuery, Snowflake information-schema,
Kinesis, GCS) and a Collector \emph{receiver} that actively polls sources are
straightforward extensions.

\section{Conclusion}
Freshness is a known and well-studied quantity that, paradoxically, has no
standard way to be \emph{emitted}. We proposed a compact OpenTelemetry semantic
convention for data staleness that unifies AoI theory with the fragmented
practice of data observability, and we released a reference implementation---a
Python SDK and a Collector processor---that is correct and cheap enough to run
everywhere. By making data age a portable, queryable telemetry signal alongside
latency and errors, we hope to let teams reason about ``how old is my data''
with the same rigor they already apply to ``how fast'' and ``how broken.''

\section*{Artifact Availability}
The semantic-convention specification, Python SDK (\texttt{otel-staleness}), Go
Collector processor (\texttt{datastaleness}), demo deployment, and benchmarks
are released as open source under Apache-2.0.

\begin{thebibliography}{99}
\footnotesize
\bibitem{kaul2012} S.~Kaul, R.~Yates, and M.~Gruteser. ``Real-time status: How
often should one update?'' In \emph{IEEE INFOCOM}, 2012.
\bibitem{yates2021survey} R.~D.~Yates et al. ``Age of Information: An
Introduction and Survey.'' \emph{IEEE J. Sel. Areas Commun.}, 2021.
arXiv:2007.08564.
\bibitem{peralta2006} V.~Peralta. ``Data Freshness and Data Accuracy: A State
of the Art.'' Technical Report, Universidad de la Rep\'ublica, 2006.
\bibitem{montecarlo} Monte Carlo Data. ``Data Freshness Explained.''
\url{https://www.montecarlodata.com/blog-data-freshness-explained/}.
\bibitem{metaplane} Metaplane. ``What is data freshness? Definition, examples,
and best practices.'' \url{https://www.metaplane.dev/blog/data-freshness-definition-examples}.
\bibitem{sifflet} Sifflet. ``What Is Data Freshness in Data Observability?''
\url{https://www.siffletdata.com/blog/data-freshness}.
\bibitem{dbtfreshness} dbt Labs. ``Source freshness.'' dbt documentation.
\url{https://docs.getdbt.com/docs/build/sources}.
\bibitem{burrow} LinkedIn. ``Burrow: Kafka Consumer Lag Checking.''
\url{https://github.com/linkedin/Burrow}.
\bibitem{otelsemconv} OpenTelemetry. ``Semantic Conventions.''
\url{https://opentelemetry.io/docs/concepts/semantic-conventions/}.
\bibitem{otelmessaging} OpenTelemetry. ``Semantic conventions for messaging
metrics.'' \url{https://opentelemetry.io/docs/specs/semconv/messaging/messaging-metrics/}.
\bibitem{otelroadmap} OpenTelemetry. ``Semantic Conventions 2026 Roadmap.''
GitHub issue, open-telemetry/semantic-conventions \#3330.
\end{thebibliography}

\end{document}
